\documentclass[11pt]{article}

\usepackage[margin=1in]{geometry}
\usepackage{amsmath,amssymb,amsthm,mathtools}
\usepackage[numbers,sort&compress]{natbib}
\usepackage[hidelinks]{hyperref}
\hypersetup{
  pdftitle={A unique cubic-coupling cancellation locus for collinear constant-vorticity capillary–gravity triads},
  pdfauthor={Xiaojian Zeng}
}
\usepackage{microtype}

\newtheorem{theorem}{Theorem}
\newtheorem{lemma}{Lemma}
\theoremstyle{remark}
\newtheorem{remark}{Remark}

\title{A unique cubic-coupling cancellation locus for collinear\\
constant-vorticity capillary--gravity triads}
\author{Xiaojian Zeng\\
\small PBC School of Finance, Tsinghua University, Beijing 100190, China\\
\small \texttt{xiaojianzeng@foxmail.com}}
\date{}

\begin{document}
\maketitle

\begin{abstract}
We classify a degeneracy of the leading resonant interaction in the
co-propagating constant-vorticity capillary--gravity wave model of Ivanov and
Martin. Let \(g,\sigma>0\), let \(k_1,k_2>0\) be distinct, and set
\(k_3=k_1+k_2\), where \(\sigma\) is surface tension divided by density.
Starting from the polynomial dispersion relations and
\(\omega_3=\omega_1+\omega_2\), we obtain all real solutions without
squaring a radical equation. There are exactly two: a common-positive
frequency solution and its simultaneous frequency--vorticity sign reversal.
On either sheet the cubic interaction numerator is a rational function of
\(q_j=k_j(g+\sigma k_j^2)\). Requiring this numerator to vanish reduces
exactly to a quartic equation for \(r=g/\sigma\). Its four linear factors are
positive and increasing on \(r\geq0\), while its value at the origin is
strictly negative. Hence every distinct positive pair \((k_1,k_2)\) has
exactly one \(r_*>0\) at which the linear resonance persists but the retained
cubic triad coefficient vanishes. The vorticity is negative on the
positive-frequency sheet and positive on its sign-reversed companion. The
result concerns only the source model's short-wave, cubic-order envelope
truncation; it neither suppresses higher-order transfer nor establishes a
finite-depth or full-Euler cancellation.
\end{abstract}

\noindent\textbf{Keywords:} capillary--gravity waves; constant vorticity;
resonant triad; Hamiltonian interaction coefficient; exact classification

\medskip
\noindent\textbf{Mathematics Subject Classification (2020):}
76B15, 76B45, 35Q35

\section{Introduction}

Three-wave resonance requires both a wave-vector relation and a frequency
relation. Constant vorticity can enable such resonances for collinear
capillary and capillary--gravity waves even when the corresponding
irrotational configuration is nonresonant
\cite{ConstantinKartashova2009,Martin2017}. Ivanov and Martin
\cite{IvanovMartin2019} derived a weakly nonlinear Hamiltonian and envelope
system for three distinct, co-propagating modes in a constant-vorticity flow.
The resulting dynamics contains a common cubic interaction numerator, so
linear resonance and nonlinear coupling are logically separate conditions.

The linear-resonance part is substantially understood. Martin
\cite{Martin2017} derived the vorticity magnitude with an explicit sign
restoration, and Martin and Basu \cite{MartinBasu2021} obtained the
equation-equivalent positive-frequency root and its uniqueness in the
capillary regime for a surface layer of piecewise-constant vorticity.
Accordingly, neither that magnitude nor negative-vorticity uniqueness is
claimed here as new. For the pure-capillary boundary \(g=0\), Constantin,
Kartashova, and Wahl{\'e}n \cite{ConstantinKartashovaWahlen2010} gave a
Hamiltonian coupling coefficient and studied resonance clusters.

Our question is instead simultaneous: when do the dispersion equations,
linear resonance, and the cubic coefficient from
\cite{IvanovMartin2019} vanish together? We first retain the unrestricted
real polynomial dispersion system and show that it has exactly two
total-sign-related sheets. We then reduce the simultaneous zero to one
strictly increasing quartic in \(r=g/\sigma\), proving existence and
uniqueness of a positive zero for every distinct positive wave-number pair.
This is not the same mechanism as the fixed-amplitude and fixed-phase
no-exchange states studied by Chabane and Choi \cite{ChabaneChoi2019}.

The scope is deliberately narrow. The calculation uses the short-wave
Fourier-multiplier replacement and the cubic, slow-envelope truncation of
\cite{IvanovMartin2019}. A zero of its cubic coefficient is not a statement
about all perturbative orders, finite depth, dissipation, or an experiment.

\section{Model and conventions}

The parent setting is a two-dimensional inviscid, incompressible fluid with a
free surface, a flat bed, gravity \(g>0\), and constant vorticity
\(\gamma=u_y-v_x\). We write
\[
  \sigma=\frac{\widetilde{\sigma}}{\rho}>0,
\]
where \(\widetilde{\sigma}\) is surface tension and \(\rho\) is density.
Thus
\[
 [k]=L^{-1},\qquad [\omega]=[\gamma]=T^{-1},\qquad
 [g]=LT^{-2},\qquad [\sigma]=L^3T^{-2}.
\]
At the short-wave level retained in \cite{IvanovMartin2019}, a positive
co-propagating wave number satisfies
\begin{equation}
  \omega^2+\gamma\omega-k(g+\sigma k^2)=0.       \label{eq:dispersion}
\end{equation}
The selected carrier frequency is the positive root
\begin{equation}
  \omega^+_\gamma(k)
  =\frac{-\gamma+\sqrt{\gamma^2+4k(g+\sigma k^2)}}{2}>0. \label{eq:positive}
\end{equation}

Let \(k_1,k_2>0\), \(k_1\neq k_2\), and \(k_3=k_1+k_2\). The resonance
condition is
\begin{equation}
  \omega_3=\omega_1+\omega_2.                   \label{eq:resonance}
\end{equation}
For convenience define
\begin{equation}
  q_j=k_j(g+\sigma k_j^2),\qquad
  C=\gamma^2+2\gamma\omega_3+2\omega_1\omega_2. \label{eq:Cdef}
\end{equation}
Direct summation of the four resonant cubic Hamiltonian terms in
\cite[Eq.~(4.31)]{IvanovMartin2019} gives \(C\). Each nonlinear right-hand
side of the three envelope equations in its Eq.~(4.32) is proportional to
\(k_j C/(2\omega_j+\gamma)\), multiplied by the appropriate quadratic
amplitude monomial. Hence \(C=0\), provided the denominators are regular,
removes precisely the displayed cubic resonant interaction.

\section{Complete real resonance sheets}

Write
\begin{equation}
\begin{gathered}
 a=q_1,\qquad b=q_2,\qquad c=q_3,\qquad
 \ell=c-a-b=3\sigma k_1k_2k_3,\\
 p=\frac{\ell}{2},\qquad
 Q=c^2-(a-b)^2,\qquad
 N=\ell^2-4ab,\qquad D=2\ell Q.                 \label{eq:invariants}
\end{gathered}
\end{equation}
Because \(q(k)=k(g+\sigma k^2)\) is strictly increasing for \(k>0\),
\(c>a,b\), and therefore
\[
 Q=(c-a+b)(c+a-b)>0,\qquad D>0.
\]

\begin{theorem}[Two-sheet classification]\label{thm:sheets}
Every real solution of the three dispersion polynomials
\(\omega_j^2+\gamma\omega_j-q_j=0\) and
\eqref{eq:resonance} is, for one
\(\varepsilon\in\{+1,-1\}\),
\begin{equation}
\begin{aligned}
 \omega_1^{(\varepsilon)}
   &=\varepsilon\sqrt{p}\,\frac{c+a-b}{\sqrt Q},&
 \omega_2^{(\varepsilon)}
   &=\varepsilon\sqrt{p}\,\frac{c-a+b}{\sqrt Q},\\
 \omega_3^{(\varepsilon)}
   &=\varepsilon\,\frac{2c\sqrt p}{\sqrt Q},&
 \gamma^{(\varepsilon)}
   &=-\varepsilon\,\frac{N}{\sqrt D}.
                                                        \label{eq:sheets}
\end{aligned}
\end{equation}
In particular, the physical positive-frequency sheet is
\(\varepsilon=+1\). The other solution is the simultaneous sign reversal
\((\omega_j,\gamma)\mapsto(-\omega_j,-\gamma)\); there are no mixed-sign
real branches.
\end{theorem}

\begin{proof}
Subtract the first two dispersion equations from the third after using
\(\omega_3=\omega_1+\omega_2\). This gives
\begin{equation}
  2\omega_1\omega_2=\ell>0.                     \label{eq:product}
\end{equation}
Thus \(\omega_1\) and \(\omega_2\) are nonzero and have the same sign, which
is also the sign of \(\omega_3\). Differences of the dispersion equations
also yield
\begin{equation}
 c+a-b=2\omega_1(\omega_3+\gamma),\qquad
 c-a+b=2\omega_2(\omega_3+\gamma).              \label{eq:ratio}
\end{equation}
The positive left-hand sides show that the common factor does not vanish.
Taking the ratio in \eqref{eq:ratio} and combining it with
\eqref{eq:product} fixes \(\omega_1,\omega_2\), up to their common sign, as
the first line of \eqref{eq:sheets}. Their sum gives the displayed
\(\omega_3\). Finally, substituting either frequency into its dispersion
equation and simplifying gives
\(\gamma=-\varepsilon N/\sqrt D\). Direct substitution of
\eqref{eq:sheets} verifies all four equations. The product, common sign,
and ratio arguments force every real solution to this list, proving
completeness.
\end{proof}

\begin{remark}
Eliminating the frequencies gives only
\[
  D\gamma^2-N^2=0,
\]
which loses the sheet label. At \(N=0\), both solutions have
\(\gamma=0\), but their three nonzero frequencies still have opposite
signs.
\end{remark}

For example, with \(g=\sigma=1\) and \((k_1,k_2,k_3)=(1,3,4)\), the two
solutions are
\begin{equation}
 (\gamma,\omega_1,\omega_2,\omega_3)
 =\frac{(-11,15,36,51)}{\sqrt{30}},\qquad
 \frac{(11,-15,-36,-51)}{\sqrt{30}}.            \label{eq:example-sheets}
\end{equation}
The first is the selected positive-frequency solution.

\section{The unique cubic-coupling zero}

\begin{lemma}\label{lem:C}
On either resonance sheet in Theorem~\ref{thm:sheets},
\begin{equation}
 C=\frac{4ab}{\ell}-\gamma^2
  =\frac{16ab(c-a)(c-b)-\ell^4}
  {2\ell\{c^2-(a-b)^2\}}.                       \label{eq:Crational}
\end{equation}
In particular, \(C\) is invariant under total sign reversal.
\end{lemma}

\begin{proof}
The dispersion equations and \eqref{eq:product} give
\[
 ab=\omega_1\omega_2(\omega_1+\gamma)(\omega_2+\gamma)
   =\frac{\ell}{2}\left(\frac{\ell}{2}
     +\gamma\omega_3+\gamma^2\right).
\]
Comparing this identity with \eqref{eq:Cdef} proves the first equality in
\eqref{eq:Crational}. Inserting \(\gamma^2=N^2/D\) and factoring proves the
second.
\end{proof}

Set \(x=k_1\), \(y=k_2\), and
\[
 r=\frac{g}{\sigma},
\]
which has dimensions \(L^{-2}\).

\begin{theorem}[Unique cancellation locus]\label{thm:zero}
For every distinct \(x,y>0\), there is exactly one \(r_*(x,y)>0\) for which
the dispersion equations, resonance, and \(C=0\) hold simultaneously.
It is the unique positive root of
\begin{equation}
\begin{split}
 P_{x,y}(r)={}&16(r+x^2)(r+y^2)
 (r+3x^2+3xy+y^2)\\
 &\quad\times(r+x^2+3xy+3y^2)
 -81x^2y^2(x+y)^4.                              \label{eq:quartic}
\end{split}
\end{equation}
Moreover, \(C<0\) for \(0\leq r<r_*\) and \(C>0\) for \(r>r_*\) on both
sheets. At \(r_*\), the positive-frequency sheet has \(\gamma<0\), while
its total-sign companion has \(\gamma>0\).
\end{theorem}

\begin{proof}
The numerator of \eqref{eq:Crational} factors exactly as
\[
 \sigma^4x^2y^2P_{x,y}(r).
\]
The removed factor and the denominator in \eqref{eq:Crational} are strictly
positive. Each of the four linear factors in the product in
\eqref{eq:quartic} is positive and strictly increasing on \(r\geq0\);
therefore \(P_{x,y}\) is strictly increasing there. At the endpoint,
\begin{equation}
 P_{x,y}(0)=-x^2y^2\{33(x+y)^4-16x^2y^2\}<0,    \label{eq:P0}
\end{equation}
and \(P_{x,y}(r)\to+\infty\) as \(r\to+\infty\). The intermediate value
theorem and strict monotonicity prove existence, uniqueness, and the sign
change.

On the positive-frequency sheet, solving \(C=0\) at the simultaneous
solution gives
\[
 \gamma=-(\omega_1+\omega_2)
       \pm\sqrt{\omega_1^2+\omega_2^2}.
\]
Both values are negative because
\(\sqrt{\omega_1^2+\omega_2^2}<\omega_1+\omega_2\). The other sheet is the
total sign reversal, so its vorticity is positive.
\end{proof}

The amplitude denominators are nonsingular at the zero:
\begin{equation}
 2\omega_j+\gamma=\sqrt{\gamma^2+4q_j}>0
 \quad\text{on the positive-frequency sheet}.   \label{eq:denominators}
\end{equation}
For \((k_1,k_2)=(1,3)\),
\[
 r_*=0.5179682428904992165\ldots,\qquad
 \gamma_*=-2.194536665566565812\ldots
\]
when \(\sigma=1\). These decimals illustrate the exact result; they are not
used in its proof.

\section{Relation to earlier formulas}

The convention in Martin \cite{Martin2017} uses
\[
 \omega=\frac{\gamma_M+\sqrt{\gamma_M^2+4q}}{2},
\]
so its scalar equation maps to the present one by
\(\gamma_M=-\gamma\). Its squared vorticity formula agrees with
\(\gamma^2=N^2/D\), while its unsquared equation restores the sign.
Likewise, after the surface-current Doppler terms cancel under
\(k_1+k_2=k_3\), the \(\lambda_1\) dispersion branch of Martin and Basu
\cite{MartinBasu2021} is \eqref{eq:positive}. Their variables reduce to
\[
 E=-\ell/\sigma,\qquad F=N/\sigma^2,\qquad
 4\sigma\Sigma^2=N^2/D.
\]
Their result therefore supplies prior linear-resonance input, but it does not
contain \(C\) or the simultaneous quartic \eqref{eq:quartic}.

The pure-capillary Hamiltonian study
\cite{ConstantinKartashovaWahlen2010} is the closest coupling-coefficient
antecedent. It lies at \(r=0\), where \eqref{eq:P0} shows that the present
coefficient is not zero; the cancellation locus is strictly inside
\(r>0\).

Two formula issues should be kept version-specific. In the inspected
arXiv-v1 rendering of \cite{IvanovMartin2019}, the displayed Eq.~(3.3) lacks
the square root needed when recovering a vorticity from the squared
resonance relation. It is dimensionally inconsistent as printed. For
\((k_1,k_2,k_3)=(1,3,4)\) and \(g=\sigma=1\), its literal value is
\(11/5760\), for which direct evaluation of the selected branch gives the
nonzero resonance residual
\[
 \omega_1+\omega_2-\omega_3=-1.3557266246176232\ldots.
\]
Equation~\eqref{eq:example-sheets} gives the substitution-consistent
replacement. The sentence after its Eq.~(4.33) also prints
\(\gamma=\omega_1+\omega_2\pm\sqrt{\omega_1^2+\omega_2^2}\). For positive
frequencies these values are positive and cannot zero the all-positive-term
expression \eqref{eq:Cdef}. The leading sign must instead be negative, as
in the proof of Theorem~\ref{thm:zero}, and dispersion and resonance must
still be imposed simultaneously.

These observations concern the checked arXiv v1 only. We did not inspect the
equation rendering in the version of record, and we do not assert the absence
of an unindexed author correction or a formal publisher correction.

\section{Physical interpretation and limitations}

At \(r_*\), the linear phase matching \eqref{eq:resonance} remains exact,
while the common numerator of all three displayed cubic envelope couplings
vanishes. Thus the leading retained resonant exchange degenerates. Calling
this \emph{no energy exchange} without an order qualifier would be too strong:
quartic and higher Hamiltonian terms, finite-depth corrections, dissipation,
or other modes may supply the first nonzero transfer. The rescaling used to
normalize a nonzero three-wave system is singular at \(C=0\) and cannot be
used there.

The analysis does not test the replacement of the finite-depth multiplier by
its short-wave sign limit. It excludes counter-propagating and noncollinear
triads, and it does not address stability, detuning width, experimental
accessibility, or the full Euler evolution. The all-negative-frequency
solution is algebraic bookkeeping for the unrestricted polynomial system,
not a second physical carrier at the same fixed vorticity.

\section{Conclusion}

The source-model resonance manifold has exactly two total-sign-related real
sheets. On that manifold, the cubic interaction coefficient admits the
closed form \eqref{eq:Crational}, and its zero is equivalent to the monotone
quartic \eqref{eq:quartic}. Every distinct positive wave-number pair
therefore determines one and only one positive gravity-to-capillarity ratio
at which linear resonance survives but the retained cubic coupling
vanishes. This exact separation of phase matching from leading nonlinear
transfer is the principal result; determining the first nonzero omitted
interaction is a separate problem.

\section*{Statements and declarations}

\paragraph{Funding.}
No funding information was supplied for this draft.

\paragraph{Competing interests.}
No competing-interest declaration was supplied for this draft.

\paragraph{Data and code availability.}
The proof is analytic and all central identities appear in the text.
Symbolic verification code, its recorded environment, and deterministic
output accompany the local source package.

\paragraph{Author contributions.}
Xiaojian Zeng is the sole prospective author. The contribution statement and
all authorship criteria require the prospective author's confirmation before
submission.

\paragraph{Generative-AI declaration.}
OpenAI Codex using the \texttt{gpt-5.6-sol} model was used on
14--15 September 2026 to assist with primary-literature comparison, prompts
for reconstructing and challenging the algebra, symbolic-verification code,
independent computational replays, and substantial manuscript drafting.
This use materially affected scientific content and presentation. The
prospective author has not yet completed final human verification or assumed
accountability for this draft; it must not be submitted until that review is
complete and the declaration is confirmed against the venue policy.

\bibliographystyle{unsrtnat}
\bibliography{references}

\end{document}
